% Reference deck: Chapter 3, Section 3.1.
\section[Introduction]{Introduction to Determinants}
\sectioncard{3.1}{Introduction to Determinants}

\begin{frame}{A Scalar Attached to a Square Matrix}
For a square matrix $A$, its \textbf{determinant} is a scalar written
\[
\det A
\qquad\text{or}\qquad
|A|.
\]
\pause
The determinant is defined only for square matrices.
\medskip

\begin{alertblock}{Notation warning}
The symbols
\[
\begin{vmatrix}a&b\\c&d\end{vmatrix}
\quad\text{and}\quad
\begin{bmatrix}a&b\\c&d\end{bmatrix}
\]
represent different objects: the first is a number, while the second is a matrix.
\end{alertblock}
\end{frame}

\begin{frame}{Determinants of Orders One and Two}
For a $1\times1$ matrix,
\[
\det\begin{bmatrix}a\end{bmatrix}=a.
\]
\pause
For a $2\times2$ matrix,
\[
\det\begin{bmatrix}a&b\\c&d\end{bmatrix}
=ad-bc.
\]
\pause
The two products follow the two diagonals:
\[
\underbrace{ad}_{\text{main diagonal}}
-
\underbrace{bc}_{\text{other diagonal}}.
\]
\end{frame}

\begin{frame}{A $2\times2$ Example}
Let
\[
A=\begin{bmatrix}4&7\\2&6\end{bmatrix}.
\]
Then
\[
\det A
=\begin{vmatrix}4&7\\2&6\end{vmatrix}
=4(6)-7(2)
=10.
\]
\pause
For a $2\times2$ matrix, the inverse formula already shows why the determinant matters:
\[
A^{-1}=\frac{1}{\det A}\begin{bmatrix}6&-7\\-2&4\end{bmatrix}.
\]
The general connection with invertibility will be developed in the next lecture.
\end{frame}

\begin{frame}{From $2\times2$ to Larger Matrices}
The $2\times2$ diagonal formula does not directly tell us how to handle an $n\times n$ matrix.
\pause

The scalable idea is to reduce the size of the problem:
\[
n\times n
\quad\longrightarrow\quad
(n-1)\times(n-1)
\quad\longrightarrow\quad\cdots\quad\longrightarrow\quad
2\times2.
\]
\pause
To do this, we delete one row and one column at a time.
\end{frame}

\begin{frame}{Submatrices and Minors}
Let $A=[a_{ij}]$ be an $n\times n$ matrix. The submatrix $A_{ij}$ is obtained by deleting row $i$ and column $j$.
\medskip

For
\[
A=\begin{bmatrix}
1&2&5&0\\
2&0&4&1\\
3&1&0&7\\
0&4&2&0
\end{bmatrix},
\]
deleting row $3$ and column $2$ gives
\[
A_{32}=\begin{bmatrix}1&5&0\\2&4&1\\0&2&0\end{bmatrix}.
\]
\pause
The \textbf{$(i,j)$-minor} is
\[
M_{ij}=\det A_{ij}.
\]
\end{frame}

\begin{frame}{Cofactors and the Sign Pattern}
The \textbf{$(i,j)$-cofactor} is
\[
C_{ij}=(-1)^{i+j}M_{ij}=(-1)^{i+j}\det A_{ij}.
\]
\pause
The signs form a checkerboard pattern:
\[
\begin{bmatrix}
+&-&+&- &\cdots\\
-&+&-&+ &\cdots\\
+&-&+&- &\cdots\\
-&+&-&+ &\cdots\\
\vdots&\vdots&\vdots&\vdots&\ddots
\end{bmatrix}.
\]
\pause
The sign belongs to the \emph{position} $(i,j)$, independent of the sign of $a_{ij}$.
\end{frame}

\begin{frame}{Recursive Definition of the Determinant}
For $n\ge2$, define the determinant by expanding across the first row:
\[
\det A
=a_{11}C_{11}+a_{12}C_{12}+\cdots+a_{1n}C_{1n}.
\]
Equivalently,
\[
\boxed{\det A=\sum_{j=1}^{n}(-1)^{1+j}a_{1j}\det A_{1j}}.
\]
\pause
For a $3\times3$ matrix,
\[
\det A=a_{11}\det A_{11}-a_{12}\det A_{12}+a_{13}\det A_{13}.
\]
\end{frame}

\begin{frame}{First-Row Expansion: Example}
Let
\[
A=\begin{bmatrix}1&5&0\\2&4&1\\0&2&0\end{bmatrix}.
\]
Expanding across the first row gives
\begin{align*}
\det A
&=1\begin{vmatrix}4&1\\2&0\end{vmatrix}
-5\begin{vmatrix}2&1\\0&0\end{vmatrix}
+0\begin{vmatrix}2&4\\0&2\end{vmatrix}\\
&=1(0-2)-5(0-0)+0(4-0)\\
&=-2.
\end{align*}
\end{frame}

\begin{frame}{Cofactor Expansion Along Any Row or Column}
\theorembox{Theorem 1}{The determinant of an $n\times n$ matrix may be computed by a cofactor expansion across any row or down any column.}
\medskip

Across row $i$:
\[
\det A=\sum_{j=1}^{n}a_{ij}C_{ij}.
\]
Down column $j$:
\[
\det A=\sum_{i=1}^{n}a_{ij}C_{ij}.
\]
\pause
Every choice gives the same determinant. The amount of arithmetic may be very different.
\end{frame}

\begin{frame}{Choose the Row or Column with Zeros}
For
\[
A=\begin{bmatrix}1&5&0\\2&4&1\\0&2&0\end{bmatrix},
\]
the third row contains two zeros. Expand there:
\begin{align*}
\det A
&=0C_{31}+2C_{32}+0C_{33}\\
&=2(-1)^{3+2}\begin{vmatrix}1&0\\2&1\end{vmatrix}\\
&=-2.
\end{align*}
\pause
Only one $2\times2$ determinant was needed.
\end{frame}

\begin{frame}{Column Expansion: Example}
Let
\[
B=\begin{bmatrix}
2&0&3\\
1&0&4\\
5&2&-1
\end{bmatrix}.
\]
The second column has two zeros, so
\begin{align*}
\det B
&=0C_{12}+0C_{22}+2C_{32}\\
&=2(-1)^{3+2}\begin{vmatrix}2&3\\1&4\end{vmatrix}\\
&=-2(8-3)=-10.
\end{align*}
\end{frame}

\begin{frame}{A Sparse $4\times4$ Determinant}
Compute
\[
\det A,
\qquad
A=\begin{bmatrix}
0&2&0&0\\
3&1&4&0\\
0&5&2&0\\
0&0&0&6
\end{bmatrix}.
\]
\pause
Expand down column $4$, then across row $1$ of the remaining matrix:
\begin{align*}
\det A
&=6\begin{vmatrix}0&2&0\\3&1&4\\0&5&2\end{vmatrix}\\
&=6\left[-2\begin{vmatrix}3&4\\0&2\end{vmatrix}\right]
=6[-2(6)]=-72.
\end{align*}
\end{frame}

\begin{frame}{Triangular Matrices}
\theorembox{Theorem 2}{If $A$ is triangular, then its determinant is the product of its main diagonal entries.}
\medskip

For example,
\[
\det\begin{bmatrix}
3&-1&4&2\\
0&5&7&1\\
0&0&-2&6\\
0&0&0&4
\end{bmatrix}
=3(5)(-2)(4)=-120.
\]
\pause
The same rule holds for lower triangular and diagonal matrices.
\end{frame}

\begin{frame}{Why the Triangular Rule Works}
Expand a triangular determinant along a row or column with many zeros.
For an upper triangular matrix,
\[
\det\begin{bmatrix}
a_{11}&*&\cdots&*\\
0&a_{22}&\cdots&*\\
\vdots&\ddots&\ddots&\vdots\\
0&\cdots&0&a_{nn}
\end{bmatrix}
=a_{11}\det\begin{bmatrix}
a_{22}&\cdots&*\\
\vdots&\ddots&\vdots\\
0&\cdots&a_{nn}
\end{bmatrix}.
\]
Repeating the expansion gives
\[
\det A=a_{11}a_{22}\cdots a_{nn}.
\]
\end{frame}

\begin{frame}{A $3\times3$ Shortcut: Sarrus' Rule}
For a $3\times3$ matrix only,
\[
\det\begin{bmatrix}a&b&c\\d&e&f\\g&h&i\end{bmatrix}
=aei+bfg+cdh-ceg-bdi-afh.
\]
\pause
For example,
\begin{align*}
\det\begin{bmatrix}1&2&3\\0&4&5\\1&0&6\end{bmatrix}
&=1(4)(6)+2(5)(1)+3(0)(0)\\
&\quad-3(4)(1)-2(0)(6)-1(5)(0)\\
&=22.
\end{align*}
\end{frame}

\begin{frame}{Shortcut Warning}
\begin{alertblock}{Sarrus' rule stops at order three}
The diagonal pattern for a $2\times2$ or $3\times3$ determinant does not generalize to $4\times4$ or larger matrices.
\end{alertblock}
\pause
For larger matrices, use one of the reliable methods:
\begin{itemize}
  \item cofactor expansion;
  \item the triangular rule when the matrix is already triangular;
  \item row reduction while tracking determinant changes, developed in Section 3.2.
\end{itemize}
\end{frame}

\begin{frame}{Practice: Choose Before You Compute}
Which expansion is most efficient?
\[
A=\begin{bmatrix}
2&0&5&0\\
1&0&3&0\\
4&7&2&6\\
0&0&1&0
\end{bmatrix}.
\]
\pause
Column $4$ contains three zeros. Expanding there leaves a $3\times3$ determinant whose second column also contains two zeros.
\pause
\begin{block}{Strategy}
Before calculating any minor, scan every row and column for zeros.
\end{block}
\end{frame}

\begin{frame}{Practice: Cofactor Sign}
Suppose an expansion uses the entry $a_{42}=-3$.
What is the signed term contributed by this entry?
\pause
\[
a_{42}C_{42}
=(-3)(-1)^{4+2}\det A_{42}
=-3\det A_{42}.
\]
\pause
The position $(4,2)$ has a positive cofactor sign. The term is negative because the matrix entry itself is negative.
\end{frame}

